\documentclass[11pt,a4paper]{article}
\usepackage[margin=2.5cm]{geometry}
\usepackage{microtype}
\usepackage{amsmath,amssymb}
\usepackage{xcolor}
\usepackage{enumitem}
\usepackage{listings}
\usepackage[hidelinks]{hyperref}

\setlength{\parskip}{0.45em}
\setlength{\parindent}{0pt}
\setlist[itemize]{leftmargin=1.4em,itemsep=0.2em,topsep=0.3em}
\lstset{
  basicstyle=\ttfamily\small,
  backgroundcolor=\color{black!4},
  frame=single,
  rulecolor=\color{black!15},
  framesep=6pt,
  xleftmargin=4pt,
  xrightmargin=4pt,
  breaklines=true,
  columns=fullflexible,
  keepspaces=true
}
\makeatletter
\def\file{\begingroup\urlstyle{tt}\def\UrlBreaks{\do\/\do\_\do\.\do\-}\Url}
\makeatother
\newcommand{\E}{\mathbb{E}}
\newcommand{\Ia}{I^{\alpha}_{0+}}

\begin{document}

{\LARGE\bfseries Riemann--Liouville integrals as expected local times\\[4pt]
on fBm zero sets}

\vspace{0.7em}
{\large Repository README}

\vspace{0.5em}
v1.0, 25 November 2026

\vspace{0.35em}
\textbf{Author:} Grok Bot, with human assistance from Giulio Prisco

\vspace{0.4em}
The full paper title is \emph{Riemann--Liouville fractional integrals as expected
local times on the zero sets of fractional Brownian motion}.
The compiled paper is \file{paper/fractpaper.pdf}.

\section*{Summary}

For $0<\alpha<1$, let $B$ be a fractional Brownian motion with Hurst index
$H=1-\alpha$, and let $\Lambda_t$ be the local time at 0 of the time-reversed
path $s\mapsto B_{t-s}$ on $[0,t]$. The paper proves that for every Borel
function $f$ with $\int_0^t |f(s)|(t-s)^{\alpha-1}\,ds<\infty$,
\[
\E\int_0^t f(s)\,\Lambda_t(ds)=\frac{\Gamma(\alpha)}{\sqrt{2\pi}}\,\Ia f(t),
\]
where $\Ia$ is the left-sided Riemann--Liouville integral. The proof uses the
occupation density formula, the Gaussian law of $B_u$, and local
nondeterminism. The identity uses fBm only through the intensity of its local
time, and is one case of a general fact: any self-similar random measure with
finite intensity gives, in the mean, a Riemann--Liouville integral whose order
is its scaling index. When the measure is an exact Hausdorff measure with a
regularly varying gauge, as for fBm zero sets, that index equals the dimension
of the set. The result is a random-fractal reading, exact in the mean, of
Nigmatullin's idea that fractional integration of order $\alpha$ is
integration over a set of dimension $\alpha$. The paper also proves $L^2$ and
almost sure convergence of the usual Gaussian-kernel approximation of local
time, with exact bias formulas, and reports Monte Carlo experiments and a
Lean~4 check of the deterministic part of the proof.

\section*{This repository}

This repository is a new version of
\url{https://github.com/giulioprisco/fracturlab/tree/main/fractalinterpretation-fbm}.
That earlier folder holds the 2025--2026 Gemini and Grok simulations, and an
AI-written preprint, that first explored the idea numerically.

The present paper is a complete rewrite. It adds a proof of the central
identity, a sharper statement (including the role of time reversal and of local
time versus Hausdorff measure), exact bias formulas, new numerics with
nonlinear test functions and a full error budget, and a partial machine check
in Lean~4. It is meant to be published under the name
\file{fractalinterpretation-fbm-2}. The intended address is
\file{https://github.com/giulioprisco/fracturlab/tree/main/fractalinterpretation-fbm-2}.

\section*{Folder map}

\begin{itemize}
\item \file{README.pdf}: this note. Its source is \file{README.tex}.
\item \file{paper/}: the compiled v1.0 paper \file{paper/fractpaper.pdf}, the
  source \file{paper/paper.tex}, and the figures in \file{paper/figs/}.
  Compile the source with \texttt{pdflatex} twice, from the \file{paper/}
  folder, so that the figures and the cross-references resolve.
\item \file{code/numerics/}: Python scripts for the new simulations, together
  with the small result files they write (\file{results/} and
  \file{pathwise/}) and the plots in \file{plots/}. The note for this folder is
  \file{code/numerics/README.pdf}. The result tables from the paper run are
  shipped as \file{results/table.pdf} and \file{results/refine_table.pdf}.
  Run the scripts from \file{code/numerics/}, because they import \texttt{fbmlt}
  from the working directory and use relative paths.
\item \file{code/rerun/}: reruns and exact checks of the two earlier
  \texttt{fracturlab} scripts. Run these from \file{code/rerun/}.
\item \file{code/requirements.txt}: Python packages used to produce the
  results.
\item \file{lean/FbmCore/}: the Lean~4 project (Lean 4.37.0-rc1 and the
  matching Mathlib release). Notes: \file{lean/README.pdf} and
  \file{lean/FbmCore/README.pdf}.
\item \file{reviews/claude-reports.pdf}: the six Claude referee reports,
  combined.
\end{itemize}

\section*{The paper}

The compiled paper is \file{paper/fractpaper.pdf}. The figures used in it are
in \file{paper/figs/}. The paper source is \file{paper/paper.tex}; compile it
with \texttt{pdflatex} twice, from the \file{paper/} folder, so that the
figures and the cross-references resolve.

\section*{Python scripts}

The results in this repository were produced with Python~3.13, NumPy~2.5.5,
SciPy~1.18.1 and Matplotlib~3.11.2. The Gemini rerun also uses the
\texttt{fbm} package (0.3.0). Install with:

\begin{lstlisting}
python -m venv .venv
source .venv/bin/activate
pip install -r code/requirements.txt
\end{lstlisting}

Unless noted otherwise, commands below are run from \file{code/numerics/}.

\subsection*{fbmlt.py}

The shared library, not a command-line script. It provides exact fBm
simulation by Davies--Harte circulant embedding, closed-form Riemann--Liouville
integrals (cross-checked against adaptive quadrature with algebraic weight to
$10^{-15}$), and exact formulae for the $\varepsilon$-smoothed and
grid-discretised expectations of the estimator.

\subsection*{run\_alpha.py}

Main Monte Carlo run for one value of $\alpha$.

\begin{lstlisting}
python run_alpha.py ALPHA [N] [LOG2N] [SEED]
\end{lstlisting}

Defaults: $N=16000$, \texttt{LOG2N}${}=15$, \texttt{SEED}${}=12345$. The
generator is initialised with \texttt{SEED}${}+1000\alpha$. The paper used
$N=40\,000$ paths, a grid of $n=2^{15}$ steps, seed \texttt{12345}, and
$\alpha\in\{0.2,0.3,0.4,0.5,0.6,0.8,0.9\}$, for example:

\begin{lstlisting}
python run_alpha.py 0.2 40000 15 12345
\end{lstlisting}

Output: \file{results/alpha_*.npz} (about 64~MB each).
These files are not included; regenerate them to rebuild the main tables and
plots from scratch.

Runtime, from the paper: five processes in parallel on an 8-core machine took
about 200~s each (about 105~s of which is fBm generation), 208~s wall time in
total, for the first five values of $\alpha$; $\alpha=0.3$ and $\alpha=0.9$
were added later with the same code and seed formula, about 170~s each.

\subsection*{analyze.py}

Reads \file{results/alpha_*.npz} and writes \file{results/table.md},
\file{results/summary.json}, and the plots in \file{plots/}. Analysis took
about 10~s. The table from the paper run is shipped as
\file{results/table.pdf} (the script still writes a Markdown table if you
rerun it). The summary and the plots are included, so you do not need the
large \texttt{.npz} files just to read the results.

\subsection*{m2\_exact.py}

Exact $\E[Y_\varepsilon^2]$ for $f\equiv 1$, by two-dimensional quadrature.
No arguments. Output (included here, about 1~s):
\file{results/m2_smoothed_exact.json}.

\subsection*{m2\_discrete.py}

Exact second moment of the discrete estimator.
\begin{lstlisting}
python m2_discrete.py 0.2 0.3 0.4 0.5 0.6 0.8 0.9
\end{lstlisting}
Writes \file{results/m2_discrete_exact.json}, which is included.

\subsection*{run\_refine.py}

Refinement run at a finer grid and a narrower kernel.

\begin{lstlisting}
python run_refine.py ALPHA N LOG2N SEED TAG
\end{lstlisting}

The paper repeated the $\alpha=0.3$ run with $n=2^{17}$ and
$\varepsilon=2.5\times10^{-4}$, and $N=200\,000$ new paths (four independent
streams of 50,000), taking about 9 minutes of wall time on four cores.
Output: \file{results/refine_<tag>.npz} (about 4~MB each,
not included).

\subsection*{analyze\_refine.py}

Reads the raw refinement dumps and writes
\file{results/refine_summary.json} and \file{results/refine_table.md}.
The table from the paper run is shipped as
\file{results/refine_table.pdf} (the script still writes a Markdown table if
you rerun it). The summary is included.

\subsection*{pathwise.py}

Pathwise test of $L^2$ convergence, with a negative control that has the same
one-point marginals but no local time.

\begin{lstlisting}
python pathwise.py ALPHA N STREAM
\end{lstlisting}

The paper used $N=20\,000$ paths per $\alpha$ and $n=2^{15}$. The random
streams were spawned from a single \texttt{SeedSequence} (entropy 20261123).
The saved runs used stream indices 0 through 6 for
$\alpha=0.2,0.3,0.4,0.5,0.6,0.8,0.9$ in that order.
Output: \file{pathwise/pw_*.npz} (about 5~MB each, not
included). The whole run, including the control, took about 20 minutes on
8 cores (about 2--5 minutes per $\alpha$).

\subsection*{pathwise\_exact.py}

Exact on-grid $\E[Y_a Y_b]$, hence the exact
$\E(Y_\varepsilon-Y_{\varepsilon/2})^2$.

\begin{lstlisting}
python pathwise_exact.py 0.2 0.3 0.4 0.5 0.6 0.8 0.9
\end{lstlisting}

Writes \file{pathwise/exact_*.json} (included). About
90~s per $\alpha$.

\subsection*{pathwise\_gridfree.py}

Grid-free value of the same mean-square difference, by nested quadrature
(reliable for $\alpha\le 0.6$, as noted in \file{code/numerics/README.pdf}).

\begin{lstlisting}
python pathwise_gridfree.py 0.2 0.3 0.4 0.5 0.6 0.8 0.9
\end{lstlisting}

Writes \file{pathwise/gridfree_*.json} (included). The
paper says the exact reference values, together with the on-grid ones, took a
few more minutes after the pathwise Monte Carlo.

\subsection*{analyze\_pathwise.py, table\_pathwise.py, plot\_pathwise.py}

\file{analyze_pathwise.py} reads \file{pathwise/pw_*.npz} and writes
\file{pathwise/summary.json}. \file{table_pathwise.py} prints the LaTeX
table rows. \file{plot_pathwise.py} writes \file{plots/pathwise.png}.
The summary and the plot from the paper run are included.

\subsection*{check\_mean.py}

Independent replicate of the $f\equiv 1$ mean test, plus low-variance
marginal checks.

\begin{lstlisting}
python check_mean.py ALPHA N STREAM
\end{lstlisting}

Included outputs: $\alpha=0.8$, $N=80\,000$, $z=-0.05$; $\alpha=0.3$,
$N=40\,000$, $z=+0.71$ (for the $\varepsilon=0.001$ mean).

\subsection*{check\_var.py}

A rough check of $\mathrm{Var}\,B_u$. It is not used in the paper.

\begin{lstlisting}
python check_var.py ALPHA N STREAM
\end{lstlisting}

\subsection*{cholesky\_check.py}

Independent exact sampler, by dense Cholesky factorisation of the fBm
covariance, with no shared code with \file{fbmlt.fbm_paths}.

\begin{lstlisting}
python cholesky_check.py ALPHA N_TOTAL BATCH [STREAM] [LOG2N]
\end{lstlisting}

\texttt{STREAM} defaults to 0 (the draft v0.5 run at $\alpha=0.8$);
\texttt{LOG2N} defaults to 12. The paper reports runs at $n=2^{12}$ with
$N=2\times10^6$ paths at $\alpha=0.2,0.3,0.9$ and $N=10^6$ at $\alpha=0.8$.
The dense covariance factor is then $4096\times4096$, about 130~MB, and costs
about 0.35~ms per path. The three later runs took about 35 minutes of compute
on 8 cores in total. The saved files used stream 0 at $\alpha=0.8$ and
streams 1, 2, 3 at $\alpha=0.2,0.3,0.9$.
Output: \file{pathwise/cholesky_*.json} (included).
At $\varepsilon=10^{-3}$ the $z$-scores are $+1.25$, $+0.20$, $+0.05$,
$-0.37$, a combined $z$ of $+0.57$.

\subsection*{cantor\_dilated.py}

No arguments. Numerical check of the randomly scaled Cantor construction in
the paper (one-dimensional quadrature of $\E M([0,u])/u^\beta$, and a Monte
Carlo estimate of $\E[(2+X)^{-\beta}]/\log 2$). It prints the values and does
not write a file.

\subsection*{Rerun of the earlier scripts}

Run these from \file{code/rerun/}.

\file{gemini_run.py} reruns the earlier Gemini script ($\alpha=0.4$,
Hosking method, 2000 steps, 1000 realizations, $\varepsilon=0.05$). It needs
the \texttt{fbm} package. Runtime: 15~min 38~s. It writes
\file{gemini_alpha0.4.png}.

\file{gemini_expectation.py} computes the exact deterministic expectation
of that estimator, including the rescale factor and the shape errors after
rescaling. No Monte Carlo.

\file{grok_run.py} reruns the earlier Grok script ($H=0.15$, $\alpha=0.85$,
exact Cholesky on a non-uniform grid). The grid size was reduced from 16,384
to 4,096, because a dense $16384^2$ Cholesky factor needed more than 8~GB.
Runtime: 44~s at $N=4096$. It writes \file{grok_H0.15.png}. The rerun gave
0.6992 against the exact 0.6850 (relative error 2.1\%).

\file{grok_closed_form_check.py} checks the closed form printed by the
Grok script. The closed form is the correct left-sided Riemann--Liouville
integral of $1+t$ at $T=1$ (1.629148 both ways). The printed mismatch (0.7230
against 0.6850) comes from comparing that left-sided integral at $T$ with the
right-sided functional at 0. Before the constant
$\Gamma(\alpha)(2\pi)^{-1/2}$ these are 1.629148 and 1.543403.

\section*{Lean files}

The project is \file{lean/FbmCore/}. It uses Lean 4.37.0-rc1 and the
matching Mathlib release. The main file, \file{FbmCore/Basic.lean}, is about
185 lines. Mathlib has no fractional Brownian motion and no local times, so
the formalization covers the deterministic core and the reduction from the
intensity of local time to the theorem, and states the fBm input as one
labelled hypothesis (\texttt{IntensityHyp}): for every Borel set $A$,
$\E L(A)=\int_{A\cap(0,t]} p_u(0)\,du$, with $p_u(0)$ the $N(0,u^{2H})$
density at 0. That hypothesis is equation (4) of the paper. The existence and
joint continuity of fBm local time, the conditional variance bound, and
uniform integrability are not formalized.

Proved, with no \texttt{sorry}:
\begin{itemize}
\item \texttt{gaussian\_at\_zero}: the $N(0,u^{2H})$ density at 0 equals
  $(2\pi)^{-1/2}u^{-H}$ for $u>0$.
\item \texttt{kernel\_reversal}: the change of variables $s=t-u$ that turns
  the intensity integral into a left-sided Riemann--Liouville kernel.
\item \texttt{integral\_power}, \texttt{normalized\_kernel\_mass},
  \texttt{expected\_mass}.
\item \texttt{det2\_mono}: the algebraic core of Step~2 of the proof (a
  2~by~2 determinant comparison).
\item \texttt{intensity\_eq}, \texttt{expected\_lintegral},
  \texttt{expected\_lintegral\_reversed}, and \texttt{main\_theorem}: from the
  intensity hypothesis to the identity for Borel $f$ with
  $\int_0^t |(t-s)^{\alpha-1}f(s)|\,ds<\infty$.
\end{itemize}

Every theorem depends only on Lean's three standard axioms
(\texttt{propext}, \texttt{Classical.choice}, \texttt{Quot.sound}), as
recorded in \file{lean/FbmCore/axioms.txt}.

\noindent
\begin{minipage}{\textwidth}
Build:
\begin{lstlisting}
cd lean/FbmCore
lake exe cache get
lake build
lake env lean scratch/Axioms.lean
\end{lstlisting}
\end{minipage}

\texttt{scratch/A.lean} is a working copy used during the formalization. A copy
kept next to the paper source during drafting was identical to
\file{FbmCore/Basic.lean} and is not duplicated here.

\section*{AI referee reports}

The paper went through six rounds of review by Claude (Anthropic). The
combined reports are in \file{reviews/claude-reports.pdf}. They cover the
earlier preprint and drafts v0.1, v0.3, v0.4, v0.5 and v0.6 (v0.2 was not
reviewed separately). The rating rose from 5/10, with a recommendation of
major revision, on the earlier preprint, to about 9/10 on v0.6, where the
recommendation was accept. These are AI reports, not independent human peer
review. The paper says so, and no human referee has yet read it.

Draft v0.6 and this repository were also reviewed by Google Gemini
(25 November 2026). That review found no mathematical issues and suggested
minor fixes: the repository visibility wording, a measurability reference
for Proposition 4.3, and a stale print line in \texttt{cantor\_dilated.py}.
All of these were applied in v1.0. Gemini also wrote one of the two earlier
exploratory scripts in the first repository.

\enlargethispage{3\baselineskip}

\section*{What is not included}

The large Monte Carlo dumps are omitted, because each file is from about 4~MB
to about 64~MB:
\begin{itemize}
\item \file{results/alpha_*.npz}, written by \texttt{run\_alpha.py}
\item \file{results/refine_*.npz}, written by \texttt{run\_refine.py}
\item \file{pathwise/pw_*.npz}, written by \texttt{pathwise.py}
\end{itemize}
Regenerate them with the commands above if you want to rebuild the tables and
plots from the raw paths. The small JSON summaries, the result tables (as
PDF), and the figures are included.

Also omitted: earlier draft files of the paper source and of the compiled PDF,
the one-off edit scripts, build logs, copies of the earlier
\texttt{fracturlab} and \texttt{fractphysics} trees, venue notes, an older
\texttt{v1\_5alpha} numerics run, and the Lean build cache (about 8~GB). After
\texttt{lake exe cache get} and \texttt{lake build}, that cache is recreated
locally.

\end{document}
